% Introduction (BGU \S 10.7.2.4).
\section{Introduction}
\label{sec:intro}

\subsection{Motivation}
Terahertz (THz) imaging is increasingly used in security screening, biomedical
diagnostics, and through-fabric inspection because the millimeter-wave band
penetrates clothing, paper, plastics, and other low-permittivity materials
that block visible light. The physics of the band exacts a price: diffraction
caps the achievable spatial resolution near the wavelength of the source,
single-pixel detectors integrate the field over a comparatively large area,
and the available noise figure is dominated by thermal background. The output
is an image that is simultaneously low resolution, low signal-to-noise, and
free of color information.

Standard deep-vision classifiers, almost all of which inherit ImageNet
pretraining, are tuned for the opposite distribution. They depend on fine
textures, sharp edges, and saturated chroma to populate the receptive fields
of the early convolutional layers. When those signals are removed the accuracy
of an off-the-shelf model collapses, and the failure mode is opaque: it is
not obvious how much of the loss is caused by reduced resolution alone, how
much by additive noise, and how much by the loss of color, nor whether
training-time interventions (regularization, augmentation, architectural
substitution) can recover any of the lost performance.

\subsection{Problem statement}
This project quantifies the accuracy degradation of three representative deep
architectures under a controlled synthetic THz-like degradation, decomposes
the loss into per-axis contributions, audits the reproducibility of the
measurements across random seeds, and documents which training-side
interventions recover lost accuracy. The architectures span the three families
that currently anchor the image-classification literature: a residual
convolutional network (ResNet50 \cite{he2016resnet}), a densely connected
convolutional network with explicit feature reuse (DenseNet121
\cite{huang2017densenet}), and a vision transformer with foveal attention
(TransNeXt-tiny \cite{shi2024transnext}).

The success criteria taken from the preliminary report are: at least 80
percent top-1 accuracy on the moderate-degradation operating point, an
identifiable architectural winner that beats the baseline by a statistically
meaningful margin, and inference latency compatible with sensor-side
deployment. Those targets are revisited in Section \ref{sec:conclusions}
against the actual measurements.

\subsection{Alternative approaches considered}
Three alternative families were reviewed during the preliminary design and
rejected in favor of synthetic-degradation training.

\textbf{(i) Denoise then classify.}\quad
A pipeline that restores the image before invoking a clean-image classifier.
The approach decouples the two problems and reuses an existing pretrained
classifier verbatim, but it requires a high-quality reference image at train
time, which a real THz sensor never produces. Restoration networks trained on
mismatched degradation distributions tend to hallucinate detail
\cite{huynh2008psnr, wang2004ssim}, which inflates PSNR / SSIM scores without
improving downstream classification.

\textbf{(ii) Super-resolution then classify.}\quad
A specialization of (i) that targets the resolution axis only. It assumes
that the dominant failure mode is spatial sampling. As Section
\ref{sec:tests} will show, resolution does dominate the per-axis attribution,
but additive noise and blur still cost between five and seven \Mpp{} each at
the moderate operating point, so a resolution-only restoration cannot recover
the full loss.

\textbf{(iii) Train-time degradation exposure.}\quad
The model is trained directly on synthetically degraded images so that the
representation adapts to the degraded distribution. This is the approach used
in modern robustness benchmarks \cite{hendrycks2021naturaladv} and the one
adopted here. It is sample-efficient (no separate restoration network), it
keeps the test-time pipeline minimal (single forward pass), and it surfaces
the architectural sensitivities directly because every backbone sees the
same degraded distribution.

\subsection{Contributions}
This work contributes:

\begin{itemize}
  \item a 402-cell training campaign across three architectures, two
        datasets, five degradation levels, and six experimental phases, all
        trained to convergence under a single fair-comparison invariant
        ($224 \times 224$ uniform input, frozen Optuna L3 hyperparameters,
        identical data pipeline);
  \item a 48-replicate multi-seed audit on the 24 L3 headline cells that
        bounds run-to-run variance below 1.5 \Mpp{};
  \item a deterministic 256-image PSNR / SSIM probe per pipeline so that
        accuracy can be compared against the physical image-quality scale,
        not only against the discrete level index;
  \item a per-axis attribution decomposition at L3 that shows resolution is
        the dominant cost ($\sim 4 \times$ blur or salt and pepper);
  \item a regularization-recovery sweep (Phase D) and a protocol simplification
        sweep (Phase B2 / B2-nr) that, together, quantify the fraction of the
        accuracy collapse that is information-bottleneck versus training noise.
\end{itemize}

\subsection{Roadmap}
Section \ref{sec:spec} reports the as-built technical specification of the
delivered software product and lists every drift from the preliminary
specification. Section \ref{sec:design} documents the engineering approach
and the mathematical model. Section \ref{sec:tests} reports the final
acceptance tests, including the per-pipeline image-quality probe and the
seven per-phase accuracy tables. Section \ref{sec:problems} enumerates the
problems encountered during implementation and how each was resolved.
Section \ref{sec:conclusions} compares the achieved performance to the
preliminary specification and discusses goal-by-goal success.
